\input{../tex-inputs/latex-leseansicht-vorspann}

\section[Arthur Schnitzler an Gustav Schwarzkopf, 24. 3. 1895]{L04118 Arthur Schnitzler an Gustav Schwarzkopf, 24. 3. 1895}
\nopagebreak\mylabel{L04118v}
\rehead{ }\normalsize\beginnumbering\briefempfaengerindex{Schwarzkopf, Gustav@\textsc{Schwarzkopf, Gustav}!zzzSchnitzler, Arthur@\emph{von Arthur Schnitzler}!1895-03-241@{24. 3. 1895}|(be}
\toendnotes[C]{\smallbreak\pagebreak[2]}
\correspDesc{Versand  durch Arthur Schnitzler am 24. 3. 1895 in Frankgasse 1
\newline{}Übermittlung  am 25. 3. 1895 in Wien 9/3 72
\newline{}Zustellung  am 25. 3. 1895 in I., Innere Stadt
\newline{}Erhalt  durch Gustav Schwarzkopf am 25. 3. 1895 in Tiefer Graben 23}\toendnotes[C]{\smallbreak}
\Standort{CUL, Schnitzler, B 96.}
\physDesc{Postkarte, 214 Zeichen
\newline{}Handschrift: Bleistift, deutsche Kurrent
\newline{}Versand: 1) Stempel: »\nobreak{}\oindex{IX., Alsergrund@\textbf{IX., Alsergrund}|pwk}Wien 9/3 72, 25. 3. 95, 6–7V\nobreak{}«.   2) Stempel: »\nobreak{}\oindex{I., Innere Stadt@\textbf{I., Innere Stadt}|pwk}Wien 1/1 1, 25. 3. 95, 8–9½V., Bestellt\nobreak{}«. }\toendnotes[C]{\smallbreak}\pstart{}{\pb}Herrn \textsc{Gustav Schwarzkopf}\pend{}\pstart{}Wien\oindex{Wien@\textbf{Wien}|pw}\pend{}\pstart{}\textsc{I. Tiefer Graben 23}\oindex{Wien@\textbf{Wien}!I., Innere Stadt@\textbf{I., Innere Stadt}!Tiefer Graben 23@\textbf{Tiefer Graben 23}|pw}.\pend{}{\bigskip}\vspace{1em}
\pstart
           \raggedleft{}{\pb}So{\geminationn}tag. –\pend
           \vspace{0.5em}
\pstart
           Lieber Freund, vergeſſen Sie nicht, \label{K_L04118-1v}\edtext{morgen Vormittag}{\lemma{\textnormal{\emph{morgen Vormittag}}}\Cendnote{\textnormal{Vgl. A. S.: \emph{Tagebuch}, 25. 3. 1895.}}}\label{K_L04118-1},
      we{\geminationn} Sie nichts beſſers vorhaben, mir das Vergnügen
      Ihres Beſuches zu ſchenken.\pend
           
\pstart
           Herzlichen Gruſs{\\[\baselineskip]}\spacefill\mbox{ArthSchnitzler}\pend
           \leftskip=0em{}\selectlanguage{ngerman}\endnumbering\briefempfaengerindex{Schwarzkopf, Gustav@\textsc{Schwarzkopf, Gustav}!zzzSchnitzler, Arthur@\emph{von Arthur Schnitzler}!1895-03-241@{24. 3. 1895}|)be}\mylabel{L04118h}
\begin{anhang}
\end{anhang}\newcommand{\dateiname}{L04118}\newcommand{\titel}{Arthur Schnitzler an Gustav Schwarzkopf, 24. 3. 1895}\newcommand{\editorInnen}{Herausgegeben von Jahnke, SelmaMüller, Martin Anton}\input{../tex-inputs/latex-leseansicht-abspann}
